\section*{Chapter \ref{ch:g12:special-relativity} --- Special Relativity: Time Dilation}

\begin{solution}{exo:g12:special-relativity:1}
$\gamma = 1/\sqrt{1 - (v/c)^2}$: $1.01$ ($0.1c$); $1.25$ ($0.6c$);
$1.67$ ($0.8c$); $10.0$ ($0.995c$).
\end{solution}

\begin{solution}{exo:g12:special-relativity:2}
$\gamma = 1.25$, so the ground measures $1.25 \times 10.0 =
\qty{12.5}{s}$. The ship's \qty{10.0}{s} is the proper time: the
metronome (one clock) is present at both beats.
\end{solution}

\begin{solution}{exo:g12:special-relativity:3}
$v = c\sqrt{1 - 1/\gamma^2}$: $\gamma = 2$ gives $v = 0.866c$;
$\gamma = 10$ gives $v = 0.995c$.
\end{solution}

\begin{solution}{exo:g12:special-relativity:4}
$v = \qty{36.1}{m/s}$, $v/c = \num{1.2e-7}$;
$\gamma - 1 \approx (\num{1.2e-7})^2/2 \approx \num{7.2e-15}$. One
second lost after $1/\num{7.2e-15} \approx \qty{1.4e14}{s}$ --- about
four million years of driving.
\end{solution}

\begin{solution}{exo:g12:special-relativity:5}
The passenger rides toward the front flash, so it reaches him first;
by the second postulate both flashes travel at the same $c$ over the
equal half-train distances \emph{in his frame}, so an earlier arrival
means an earlier strike. Simultaneity is frame-dependent: each verdict
is correct in its own frame.
\end{solution}

\begin{solution}{exo:g12:special-relativity:6}
$\gamma = 1/\sqrt{1 - 0.99^2} = 7.09$; lab lifetime
$7.09 \times \num{2.6e-8} = \qty{1.8e-7}{s}$. Classical range
$0.99 \times \num{3.00e8} \times \num{2.6e-8} \approx \qty{7.7}{m}$;
actual $7.09 \times 7.7 \approx \qty{55}{m}$.
\end{solution}

\begin{solution}{exo:g12:special-relativity:7}
$\gamma = 1.5$ at $v = c\sqrt{1 - 1/2.25} \approx 0.75c$;
$\gamma = 3$ at $\approx 0.94c$. The asymptote: $\gamma$ (and the
energy cost of further speed) diverges as $v \to c$ --- no finite
effort brings a massive body to light speed.
\end{solution}

\begin{solution}{exo:g12:special-relativity:8}
$\Delta t_0 = 2L/c = 3.0/\num{3.00e8} = \qty{1.0e-8}{s}$; at $0.8c$,
$\gamma = 5/3$, so $\Delta t = \qty{1.7e-8}{s}$.
\end{solution}

\begin{solution}{exo:g12:special-relativity:9}
Proper time belongs to the clock present at both events: (a) the
wearer; (b) the muon; (c) the pilot --- only the ship is at both the
departure and the arrival.
\end{solution}

\begin{solution}{exo:g12:special-relativity:10}
$\gamma - 1 \approx 7700^2/(2 \times \num{9.0e16}) \approx
\num{3.3e-10}$; over \qty{1.58e7}{s} the astronaut trails by
$\num{3.3e-10} \times \num{1.58e7} \approx \qty{5.2e-3}{s}$ --- about
five milliseconds younger.
\end{solution}

\begin{solution}{exo:g12:special-relativity:11}
$E = \num{1.0e-3} \times (\num{3.00e8})^2 = \qty{9.0e13}{J}$. A
\qty{1.0}{GW} station delivers $\num{1.0e9} \times 86400 \approx
\qty{8.6e13}{J}$ per day --- one gram is one day of its output, and
$\num{9.0e13}/\num{8.5e5} \approx \num{e8}$: a hundred million
speeding cars.
\end{solution}

\begin{solution}{exo:g12:special-relativity:12}
Square and invert: $\gamma^2(1 - v^2/c^2) = 1$, so
$v^2/c^2 = 1 - 1/\gamma^2$, hence $v = c\sqrt{1 - 1/\gamma^2}$. For
$\gamma = 30$: $v = c\sqrt{1 - 1/900} = 0.99944c$.
\end{solution}

\begin{solution}{exo:g12:special-relativity:13}
Contracted thickness $L = \qty{15}{km}/30 = \qty{500}{m}$; sweep time
$500/(0.99944 \times \num{3.00e8}) \approx \qty{1.7e-6}{s} <
\qty{2.2}{\micro s}$. Both frames agree the muon reaches the ground ---
one blames a stretched lifetime, the other a shrunken atmosphere.
\end{solution}

\begin{solution}{exo:g12:special-relativity:14}
Lab: $3200/\num{3.00e8} \approx \qty{1.1e-5}{s}$ (speed
$\approx c$). Electron frame: tube contracted to
$3200/\num{1.0e5} = \qty{3.2}{cm}$, crossed in
$\num{1.1e-5}/\num{1.0e5} \approx \qty{1.1e-10}{s}$.
\end{solution}

\begin{solution}{exo:g12:special-relativity:15}
Earth: $4.2/0.9 \approx \qty{4.7}{years}$. With
$\gamma = 1/\sqrt{1 - 0.81} = 2.29$, the probe's clock logs
$4.7/2.29 \approx \qty{2.0}{years}$.
\end{solution}

\begin{solution}{pb:g12:special-relativity:1}
\textbf{1.} $c\,\Delta t_0 = \num{3.00e8} \times \num{2.2e-6} \approx
\qty{660}{m}$.

\textbf{2.} $\num{15000}/660 \approx 23$ lifetimes.

\textbf{3.} $\mathrm{e}^{-23} \approx \num{e-10}$: fewer than one muon
in ten billion should arrive, yet the flux is one per \unit{cm^2} per
minute --- classical physics is flatly contradicted by the sky.

\textbf{4.} Ground lifetime $30 \times \qty{2.2}{\micro s} =
\qty{66}{\micro s}$; range $\approx \num{3.00e8} \times \num{6.6e-5}
\approx \qty{20}{km} > \qty{15}{km}$: the arrival is explained.

\textbf{5.} $v = c\sqrt{1 - 1/900} = 0.99944c$.

\textbf{6.} $\Delta t_0 = 2L/c = 3.0/\num{3.00e8} = \qty{1.0e-8}{s}$.

\textbf{7.} Horizontal $v\,\Delta t/2$; vertical $L$; the slanted
half-path measures $c\,\Delta t/2$ because the second postulate fixes
the photon's speed at $c$ in the ground frame too.

\textbf{8.} $(c\,\Delta t/2)^2 = L^2 + (v\,\Delta t/2)^2$ gives
$\Delta t^2(c^2 - v^2) = 4L^2$, so $\Delta t =
(2L/c)/\sqrt{1 - v^2/c^2} = \gamma\,\Delta t_0$.

\textbf{9.} $\gamma = 1/\sqrt{1 - 0.36} = 1.25$;
$\Delta t = \qty{1.25e-8}{s}$.

\textbf{10.} The ship's: its one clock is present at both ticks, so it
reads the proper time; the ground needs two synchronized clocks.

\textbf{11.} $\gamma = 1/\sqrt{1 - 0.64} = 1/0.6 = 5/3 \approx 1.67$.

\textbf{12.} $10.0/(5/3) = \qty{6.0}{years}$.

\textbf{13.} Outbound $\qty{5.0}{years}$ at $0.8c$: $4.0$ light-years.

\textbf{14.} $10.0 - 6.0 = \qty{4.0}{years}$: the astronaut is four
years younger than her twin.

\textbf{15.} The situation is not symmetric: only the astronaut turns
around --- she accelerates and changes inertial frame, while the Earth
twin stays in one --- so only her clock logs the shorter time.

\textbf{16.} $v = 2\pi \times \num{2.66e7}/\num{43200} \approx
\qty{3.9e3}{m/s}$.

\textbf{17.} $\gamma - 1 \approx (\num{3.87e3})^2/(2 \times
\num{9.0e16}) \approx \num{8.3e-11}$.

\textbf{18.} $\num{8.3e-11} \times \num{86400} \approx \qty{7.2e-6}{s}
= \qty{7.2}{\micro s}$ per day slow.

\textbf{19.} $c \times \num{7.2e-6} \approx \qty{2.2}{km}$ of position
error per day.

\textbf{20.} Net drift $45 - 7.2 \approx \qty{38}{\micro s}$ per day
fast; error $c \times \num{38e-6} \approx \qty{11}{km}$ per day ---
the relativity bill that every GPS clock pays in advance, detuned on
the ground so it ticks true in orbit.
\end{solution}
